\section*{第 \ref{ch:b3:electromagnetism-in-matter} 章——物质中的电磁学}

\begin{solution}{exo:b3:electromagnetism-in-matter:1}
(a) 面上电荷为 $+P$，位于 $\vect P$ 所指的一侧，另一侧为 $-P$。(b)
两块无限大带电面，面电荷为 $\pm P$，其间产生 $P/\varepsilon_0$，
方向由 $+$ 指向 $-$，即与 $\vect P$ \emph{相反}。(c)
$\vect D = \varepsilon_0\vect E + \vect P = 0$——这是必然的：
无自由电荷，且 $\vect D$ 的法向分量
与外侧真空连续。(d) 条形磁铁：偶极矩密度被冻结，
其场由自身的束缚源维持。
\end{solution}

\begin{solution}{exo:b3:electromagnetism-in-matter:2}
(a) $C = \varepsilon_{\text{r}}C_0 = \qty{400}{pF}$。(b) $Q =
\qty{1.2}{nC}$ 保持不变：$V = Q/C = \qty{3}{V}$。(c)
$\tfrac12C_0V_0^2 = \qty{7.2}{nJ}$ 为插入前，$\qty{1.8}{nJ}$ 为插入后：
缺失的 \qty{5.4}{nJ} 转化为机械功——边缘
场\emph{把介质板拉入}。(d) 接上电池后，$V$ 保持
\qty{12}{V} 不变，$Q$ 变为四倍，能量升至 \qty{28.8}{nJ}——
电池既供给新增的场能，也供给拉入介质板的功。
\end{solution}

\begin{solution}{exo:b3:electromagnetism-in-matter:3}
水：抗磁性——电子全部成对，仅有拉莫尔感应
引起的反向作用。铝：顺磁性（传导电子的自旋）。
液氧：强顺磁性——O$_2$ 含有\emph{两
个未成对}电子（其反键轨道按洪特规则填充，
见\cref{ch:b3:identical-particles}）；它会明显地
\omterm{prop:b3:grand-canonical:langmuir}{吸附}在两磁极之间。镍：铁磁性。磁铁明显
吸引镍和液氧，并极轻微地排斥
水——即强磁铁在水面上造成的凹陷。
\end{solution}

\begin{solution}{exo:b3:electromagnetism-in-matter:4}
(a) $H = nI = \qty{3000}{A/m}$。(b) $M = B/\mu_0 - H =
\qty{1.27e6}{A/m}$。(c) $B/\mu_0H \approx 420$。(d)
等效束缚面电流为 $M \approx \qty{1.27e6}{A}$ 每
米——是线圈自身 $nI = \qty{3000}{A/m}$ 的四百倍：
铁芯几乎承担了全部工作。
\end{solution}

\begin{solution}{exo:b3:electromagnetism-in-matter:5}
(a) 取圆柱形高斯面，其内只包围自由电荷：
$D = \lambda/2\pi r$，沿径向。(b) $E =
\lambda/2\pi\varepsilon_0\varepsilon_{\text{r}}r$，$P =
(\varepsilon_{\text{r}}-1)\lambda/2\pi\varepsilon_{\text{r}}r$。
(c) $\sigma_{\text{b}}(a) = -P(a)$：单位长度上的
$-(\varepsilon_{\text{r}}-1)\lambda/\varepsilon_{\text{r}}$；
在 $r = b$ 处，$+P(b)$：单位长度上的
$+(\varepsilon_{\text{r}}-1)\lambda/\varepsilon_{\text{r}}$ ---
总和为零，\omterm{def:b3:electromagnetism-in-matter:polarisation}{束缚电荷}必然如此。(d) $V = \int E\,\dd r =
(\lambda/2\pi\varepsilon_0\varepsilon_{\text{r}})\ln(b/a)$，因此
$C/\ell = 2\pi\varepsilon_0\varepsilon_{\text{r}}/\ln(b/a)$：
乘以 $\varepsilon_{\text{r}}$。
\end{solution}

\begin{solution}{exo:b3:electromagnetism-in-matter:6}
(a) $n = P/k_{\text{B}}T = \qty{2.4e25}{m^{-3}}$，因此
$\alpha/\varepsilon_0 = \chi/n \approx \qty{2.2e-29}{m^3}$ ---
半径为 $\approx\qty{1.7}{\angstrom}$ 的球体的体积：
\omterm{prop:b3:electromagnetism-in-matter:mossotti}{极化率}\emph{即}原子体积，以 $\varepsilon_0$ 为单位。
(b) $n\alpha/3\varepsilon_0 = 0.15$，因此 $(\varepsilon_{\text{r}} -
1)/(\varepsilon_{\text{r}} + 2) = 0.15$，可得
$\varepsilon_{\text{r}} = 1.53$——与实测值完全吻合。
(c) 此处局域场修正是 15\,\% 的效应；氩为
非极性，因此不存在可能失控放大的取向贡献。(d)
极性分子（水、HCl）：永久偶极子占主导并
相互作用；朗之万–德拜（对稠密液体还包括翁萨格）理论
取代了简单的空腔论证。
\end{solution}

\begin{solution}{exo:b3:electromagnetism-in-matter:7}
(a) $\chi_{\text{e}} = np^2/3\varepsilon_0k_{\text{B}}T \approx
11.5$。(b) $E = k_{\text{B}}T/p \approx \qty{6.7e8}{V/m}$——
十倍于击穿场强：水始终深处于
线性区。(c) 在 \qty{2.45}{GHz} 下，偶极子几乎能跟上
但有滞后：响应的异相部分在每个周期对水做净功，
即介质加热。在冰中，分子
被锁定在晶格上；其转动弛豫位于
千赫兹，因此在吉赫兹下冰几乎不吸收（除霜循环
以轻柔的脉冲进行，使融化的水完成加热）。(d) 水：
为负（$1/T$ 取向响应）；氩：基本为零——
感应\omterm{prop:b3:electromagnetism-in-matter:mossotti}{极化率}与温度无关。
\end{solution}

\begin{solution}{exo:b3:electromagnetism-in-matter:8}
(a) 拉莫尔频移使环流改变量为 $\Delta I
= -e\Delta\omega/2\pi = -e^2B/4\pi m_{\text{e}}$，因此
$\Delta\mu = \Delta I\cdot\pi r^2 = -e^2r^2B/4m_{\text{e}}$，
与 $\vect B$ 相反，无论电子朝哪个方向绕行——
楞次定律在原子尺度上的体现。(b) 代入给定数值，
$\chi_{\text{m}} \approx -9.5\times10^{-6}$：与实测
值 $-9\times10^{-6}$ 相符，基于环形模型。(c) 感应磁矩
来自轨道畸变，而不是玻尔兹曼式的
取向与无序之间的竞争——式中完全不含 $T$。(d) 水是
抗磁性的，因此青蛙被推向弱场区；力
密度 $\propto \chi\,\nabla(B^2/2\mu_0)$ 在均匀场中为零
场中为零——悬浮需要 $B\,\dd B/\dd z$ 足够大，以
平衡 $\rho g$，因此需要 \qty{16}{T} 的孔径。
\end{solution}

\begin{solution}{exo:b3:electromagnetism-in-matter:9}
(a) $\chi_{\text{m}} = \mu_0n\mu_{\text{B}}^2/3k_{\text{B}}T
\approx 1.7\times10^{-5}$（\qty{300}{K}），$5\times10^{-3}$（
\qty{1}{K}）。(b) 外推可得 $\chi \to 1$ 位于 \qty{5}{mK} 附近——
但偶极耦合与交换耦合会使真实盐类率先
有序（或冻结）：居里定律是高温定律。(c) 在
\qty{1}{K} 下磁化后，自旋熵被挤入热库；
若系统孤立，$S$ 固定，降低 $B$ 使自旋
各自回收 $k_{\text{B}}\ln 2$——能量来自
晶格，其温度下降（实际中可达毫开尔文）。(d)
铁磁体的自旋在 $T_{\text{c}}$ 以下自发有序：其
熵不再由磁场控制，且\omterm{def:b3:electromagnetism-in-matter:hysteresis}{磁滞}
会使能量耗散，而非降温。
\end{solution}

\begin{solution}{exo:b3:electromagnetism-in-matter:10}
(a) $NI = B(\ell/\mu_0\mu_{\text{r}} + e/\mu_0)$：其中
$\ell/\mu_{\text{r}} = \qty{0.15}{mm}$ 等效，相对于 $e =
\qty{4}{mm}$，$NI \approx 3960$ 安培匝，$I \approx
\qty{7.9}{A}$。(b) $4/4.15 \approx 96\,\%$ 的磁动势穿过
气隙。(c) $\Phi = NI/(\mathcal R_{\text{铁}} +
\mathcal R_{\text{隙}})$，其中 $\mathcal R = \ell/\mu A$：
磁动势 $NI$ 相当于电压，磁通量相当于电流，
磁阻相当于电阻——此处 $\mathcal R_{\text{隙}}
\approx 2.0\times10^{6}\,\unit{H^{-1}} \approx
27\,\mathcal R_{\text{铁}}$。(d) 超过 \qty{1.8}{T} 后，铁的
增量 $\mu_{\text{r}}$ 骤降趋近 1，其磁阻剧增，
额外电流几乎换不来额外磁通，若将电动机推至此处，
转矩不再增加，而绕组却会过热烧毁。
\end{solution}

\begin{solution}{exo:b3:electromagnetism-in-matter:11}
(a) $P = wfV = 40\times50\times\num{4e-3} = \qty{8}{W}$ ---
可接受。(b) $\qty{1.2}{kW}$：铁芯会发红；硬钢
因自身特性而不宜用于交流。(c) 涡流电动势
与回路面积及 $f$ 成正比；叠片（或粉末化，或采用
绝缘铁氧体）切断涡流回路，削减 $f^2$ 损耗——
这正是铁氧体主导兆赫频段的原因。(d) 存储器要求
存储位的势垒 $\sim\mu_0H_{\text{c}}M_{\text{s}}V \gg
k_{\text{B}}T$ 在十年间成立（超顺磁极限），即
$H_{\text{c}}$ 尽可能大（以便可写入）；变压器则要求薄环的
$H_{\text{c}} \to 0$。同一种材料无法同时兼顾两端：
电机用软磁材料，存储器用硬磁材料。
\end{solution}

\begin{solution}{exo:b3:electromagnetism-in-matter:12}
(a) 沿轴向磁化的长针是束缚电流的长螺线管：内部
$H \approx 0$（全部 $B$ 来自 $M$）。薄板沿板面磁化，
是可能最短、最粗的螺线管：其内部回流场
给出 $H = -M$。(b)
短粗磁体工作于 $H = -N_{\text{d}}M$ 处，深入其
退磁曲线：畴壁运动会逐步蚕食\omterm{def:b3:electromagnetism-in-matter:hysteresis}{剩磁}；针状磁体
位于 $H \approx 0$ 处，并保留其\omterm{def:b3:electromagnetism-in-matter:hysteresis}{剩磁}。(c) 因此指南针采用
\emph{磁针}，以及磁带与早期硬盘中细长的单畴颗粒
——形状各向异性充当存储的
保险。(d) 磁通守恒与 $\oint\vect H\cdot
\dd\vect\ell = 0$ 给出球内 $H_{\text{内}} =
H_0 - M/3$；当 $\mu_{\text{r}}\to\infty$ 时，铁使
$H_{\text{内}}\to0$，于是 $M \to 3H_0 = 3B_0/\mu_0$：球体
仅能将磁场集中至三倍，无论铁多么“优良”。
\end{solution}

\begin{solution}{pb:b3:electromagnetism-in-matter:1}
\textbf{1.} $B \approx \mu_0NI/\ell = \qty{13}{mT}$——冰箱磁铁
量级的磁场，源自 \qty{250}{W} 的线圈。
\textbf{2.} 线圈的微小 $H$ 使\omterm{def:b3:electromagnetism-in-matter:hysteresis}{磁畴}解钉并转动，其束缚电流
随即与线圈电流同步环流：铁自身额外贡献
$\mu_0M \gg \mu_0H$。
\textbf{3.} $\mu_0M_{\text{s}} = \qty{2.1}{T}$：无论线圈多强，
也无法从铁的磁极中再拉出更多。
\textbf{4.} $M = 2.2\mu_{\text{B}}n = 2.2\times\num{9.27e-24}
\times\num{8.5e28} \approx \qty{1.7e6}{A/m}$——相符。
\textbf{5.} 任何锈迹或砂粒都相当于额外串联的气隙，正处于
磁阻最关键的位置；十分之一毫米的水垢
会明显吞噬吸力。
\textbf{6.} 磁通离开中心北极，穿过接触间隙进入钢卷，
沿钢材流动，并
经外侧环形南极及磁轭返回——负载
闭合了磁路。
\textbf{7.} $NI = H_{\text{铁}}\ell + 2H_{\text{隙}}e =
B(\ell/\mu_0\mu_{\text{r}} + 2e/\mu_0)$.
\textbf{8.} 单位面积上：$\ell/\mu_{\text{r}} = \qty{0.6}{mm}$
与 $2e = \qty{3}{mm}$ 相比，三毫米空气占据
$\sim83\,\%$ 的安培匝。
\textbf{9.} 线性公式给出 $B = \mu_0NI/(\ell/
\mu_{\text{r}} + 2e) \approx \qty{3.5}{T}$——\emph{高于}
第 3 题的上限：磁轭饱和，实际给出 $B \approx
\qty{1.7}{T}$，尽管线圈已不堪重负。
\textbf{10.} $B^2/2\mu_0 \approx \qty{1.15e6}{Pa} \approx
\qty{1150}{kPa}$——相当于十一个大气压的拉力。
\textbf{11.} $F = \num{1.15e6}\times0.12 \approx \qty{140}{kN}
\approx 14$ 吨。
\textbf{12.} 可以：14 吨对 4 吨的要求
（2 吨 $\times$ 安全系数 2）——余量正是为
第 13 题而设。
\textbf{13.} 取 $e = \qty{4}{mm}$ 且 $A_{\text{有效}} =
\qty{0.04}{m^2}$：$B \approx \qty{1.5}{T}$，压强
$\approx\qty{870}{kPa}$，力 $\approx\qty{35}{kN} \approx 3.5$
吨——十二吨的“平板”额定值
仅勉强满足实际废钢的规格，因此承载能力总是按
打磨平板的条件标注。
\textbf{14.} \omterm{def:b3:electromagnetism-in-matter:hysteresis}{剩磁}：断电后，铁停留在其回线上的
$B_{\text{r}}$ 处——磁轭此时成为弱永磁体，
轻载仍可悬挂其上。
\textbf{15.} 软铁的细回线使 $B_{\text{r}}$ 与
$H_{\text{c}}$ 都很小：磁体必须按指令\emph{遗忘}磁性；
硬钢则会使起重器变成永磁体，
却带着一个毫无作用的开关。
\textbf{16.} 施加经标定的反向脉冲，使材料达到
$-H_{\text{c}}$（回线的零 $B$ 交点），并释放
负载。
\textbf{17.} 衰减的交流扫描描绘出越来越小的嵌套回线，
螺旋趋向原点：即退磁状态——与船体及旧式阴极射线管屏幕所用的
消磁方法相同。
\textbf{18.} \qty{800}{\celsius} 即 \qty{1073}{K} $> T_{\text{c}}$：
板坯呈\emph{顺磁性}——$\mu_{\text{r}} \approx 1$，
磁路断开，吸力骤降；热轧厂搬运板坯
须用夹钳，而不用磁铁。
\textbf{19.} 将 $e$ 加倍使 $B$ 减半、力变为四分之一
（气隙主导的磁路）；将 $NI$ 加倍可同时恢复两者——但
$P = RI^2$ 增至四倍：气隙的代价以铜耗热量支付。
\textbf{20.} $R = \rho L/A = \num{1.7e-8}\times1000/
\num{1.6e-5} \approx \qty{1.1}{\ohm}$.
\textbf{21.} $P = RI^2 \approx \qty{660}{W}$；按
$8\,\text{小时}\times0.6$：$\approx\qty{3.2}{kWh}$ 每班——
移动数百吨物料只需几欧元电费。
\textbf{22.} 气隙体积 $2Ae = \qty{3.6e-4}{m^3}$，其能量密度为
\qty{1.15e6}{J/m^3}：$\approx\qty{410}{J}$——不足
一秒的线圈发热量。稳态下，几乎全部电能输入都
转化为铜耗热；磁场一经建立，其代价
仅为维持电流。
\textbf{23.} 断开电路使 $LI$ 骤降为零：
线圈会产生巨大的感应尖峰，使触点燃弧；
续流二极管（或电阻）使储存的磁场能量悄然
化为热量耗散——而回线面积每周期都在铁中耗散，
不论如何皆然。
\textbf{24.} 提升力需要电流：电流一断，钢卷即落下。
\omterm{def:b3:electromagnetism-in-matter:hysteresis}{剩磁}仅保留象征性的吸力——以千克计，而非吨。
故须设立通道规定，并对任何悬于人上方的吊物，
采用电池备用（或永磁加释放线圈）的起重器。
\textbf{25.} 上限：$\mu_0M_{\text{s}} \approx \qty{2}{T}$。
吸力：$B^2A/2\mu_0$，由毫米量级的气隙支配。控制：
软铁、小 $B_{\text{r}}$，以消磁脉冲使其释放。成本：
$RI^2$，相当于千瓦级的加热器——本章全部内容都挂在
起重机上。
\end{solution}
